\section{Dataset and Problem Formulation}
\label{sec:dataset}

\subsection{MQTT-IoT-IDS2020 as released}

MQTT-IoT-IDS2020~\cite{hindy2020mqttdataset} was generated on a simulated
network of twelve MQTT sensors, a broker, a simulated camera and an attacker,
across five capture scenarios: normal operation, aggressive scan, UDP scan,
Sparta SSH brute-force and MQTT brute-force. The release provides raw packet
captures together with three derived abstractions---packet features,
unidirectional-flow features and bidirectional-flow features. This study uses
the two flow abstractions, which is the form in which the dataset is intended
for machine learning.

Table~\ref{tab:dataset} reports what the release actually contains, measured
rather than quoted. The uni-flow abstraction holds 495{,}525 records over
nineteen columns and the bi-flow abstraction 259{,}379 records over
thirty-two. Across both, traffic involves twenty-four distinct IP endpoints
and, in the uni-flow case, one hundred and six ordered endpoint pairs.

We assigned each endpoint a role from its observed traffic rather than from the
testbed description, using rules keyed on the share of inbound flows arriving
on a given port, outbound fan-out, and the number of distinct attack classes an
endpoint sources. The rules recover exactly twelve MQTT sensors, one broker,
one attacker, one SSH target, one gateway, one receive-only external host and
seven lower-volume hosts that appear chiefly in the scan captures. That the
sensor count matches the published testbed description independently supports
the assignment. Figure~\ref{fig:testbed} shows the result and the per-endpoint
evidence accompanies the code.

One property of the labelling deserves emphasis because it affects
interpretation throughout: the dataset marks \emph{both} directions of an
attack exchange. The broker's replies to a brute-force attempt carry
\texttt{is\_attack}~$=1$, and during the aggressive scan the scanned sensors'
own responses are likewise attack-labelled. An endpoint can therefore source
attack-labelled flows without being the attacker, which is why our attacker
rule keys on the number of \emph{distinct} attack classes sourced---a condition
only the address 192.168.2.5 satisfies.

\subsection{Two construction properties that constrain the experiment}

\subsubsection{Attack captures contain benign background traffic}

Each attack capture also records the testbed's ordinary traffic. The uni-flow
Sparta capture, for instance, contains 154{,}175 flows with
\texttt{is\_attack}~$=0$ alongside 28{,}232 attack flows; across the uni-flow
release, 204{,}350 records---41.2\% of the total---are benign flows captured
during an attack scenario. Deriving the multi-class label from the source
capture alone, as is natural when the captures are named after the attacks,
would relabel all of them as attacks and destroy the meaning of the task.

We therefore construct the multi-class label from both available sources, each
doing the job it is fit for: \texttt{is\_attack} determines \emph{whether} a
flow is an attack, and the source capture determines \emph{which} attack it is.
Formally, for flow $f$ with capture $c(f)$,
\begin{equation}
y_{\mathrm{mc}}(f) =
\begin{cases}
0 & \text{if } \texttt{is\_attack}(f) = 0,\\
\kappa(c(f)) & \text{otherwise,}
\end{cases}
\label{eq:label}
\end{equation}
where $\kappa$ maps a capture to its attack class in
$\{1,\dots,4\}$. The attack class is never inferred from \texttt{is\_attack}
alone, and the capture identity never overrides a benign flag. The capture is
retained per record so that provenance stays auditable. The binary label is the
dataset's own \texttt{is\_attack} column unchanged.

Under Eq.~\eqref{eq:label} the uni-flow class distribution is 376{,}186 normal,
28{,}874 MQTT brute-force, 39{,}797 aggressive scan, 22{,}436 UDP scan and
28{,}232 Sparta SSH, giving a minority class at 4.5\% of the data.

\subsubsection{Attack classes originate from one or two fixed addresses}

Every attack class in the release is sourced by at most three addresses, and
the attacker 192.168.2.5 sources all four. Victims are better distributed:
aggressive scan reaches eighteen distinct destinations and UDP scan eighteen,
whereas MQTT brute-force reaches four and Sparta SSH two. This asymmetry
decides which endpoint-holdout protocols are well posed, and we return to it in
Section~\ref{sec:methodology}.

\subsubsection{Consequence for the base-paper comparison}

The base paper~\cite{norouzi2026reliable} states that the MQTT-IoT-IDS2020
dataset ``consists of 19 features and 33{,}079 instances''. The uni-flow
abstraction has exactly nineteen columns, and the uni-flow MQTT brute-force
capture has exactly 33{,}079 records. The reported evaluation therefore appears
to cover one capture of the release rather than the release, whose uni-flow
abstraction contains 495{,}525 records over five captures. We draw no inference
about the method from this; we note only that the evaluated populations differ,
so the 99.997\% figure and any figure computed here describe different
problems. We consequently quote it as a literature value with its protocol
attached and never as a measurement produced by this study, and we additionally
report our own reproduction of GA-EBT under our protocols.

\subsection{No temporal information is available}

Neither flow abstraction contains a timestamp column; we verified this
programmatically over all ten CSV files rather than assuming it. Row order is
consequently not treated as time, no artificial timestamps are constructed, and
no temporal claim appears anywhere in this paper. A temporal extension is
possible from the released packet captures and is stated as future work in
Section~\ref{sec:conclusion}.

\subsection{Problem statement}

Let $\mathcal{F} = \{f_1,\dots,f_M\}$ be the set of flow records, each with a
source endpoint $s(f)$, a destination endpoint $d(f)$ and an attribute vector
$z_f \in \mathbb{R}^{D}$ comprising the released traffic features and the
deterministic derivations of Section~\ref{sec:methodology}. The task is to
learn $g: f \mapsto y$ for $y \in \{0,1\}$ (binary) or
$y \in \{0,\dots,4\}$ (multi-class), with per-flow granularity preserved: the
unit of prediction is the individual flow record, not an endpoint, a pair or a
time window. The endpoint identifiers $s(f)$ and $d(f)$ are available to the
model as graph structure but never as numeric features; \texttt{is\_attack} is
never an input.
